\chapter{Manual de Referencia Técnica MDA-Audio-CIM-Relations-Machines}

Este documento especifica el lenguaje \textbf{MDA-Audio-CIM-Relations-Machines} de relaciones entre máquinas definidas en el documento de especificación de máquinas CIM. Este documento trabaja al mismo nivel de abstracción (CIM), orquestando la comunicación de alto nivel entre máquinas independientes a escala de proyecto.

En el CIM no se tiene porqué hablar de implementaciones concretas, sino de \textbf{elementos} abstractos que capturan la intención del diseño de síntesis de la máquina.

\section{Estructura del Modelo}

El documento de relaciones define cómo fluye la información entre máquinas mediante una colección de vínculos abstractos.

\subsection{Modelo Raíz (\texttt{Document})}

El objeto raíz contiene la descripción global del sistema de interconexiones.

\begin{table}[H]
\centering
\small
\begin{tabular}{@{}llp{5cm}l@{}}
\toprule
Atributo & Tipo & Descripción & Restricciones \\ \midrule
\texttt{description} & STRING & Propósito general de las interconexiones. & Opcional. 1 a 1000 car. \\
\texttt{relations} & ARRAY[Relation] & Lista de objetos de relación entre máquinas. & Obligatorio. \\ \bottomrule
\end{tabular}
\end{table}

\textbf{Ejemplo de Documento:}
\begin{lstlisting}[language=json]
{
  "description": "Simple Lead to Mixer routing.",
  "relations": [ ... ]
}
\end{lstlisting}

\subsection{Objeto de Relación (\texttt{Relation})}

Define un vínculo direccional entre dos máquinas CIM independientes.

\begin{table}[H]
\centering
\small
\begin{tabular}{@{}llp{4cm}l@{}}
\toprule
Atributo & Tipo & Descripción & Restricciones \\ \midrule
\texttt{id} & STRING & Identificador único de la relación. & Obligatorio. 36 car. \\
\texttt{source} & STRING & ID de la máquina de origen. & Obligatorio. 36 car. \\
\texttt{destination} & STRING & ID de la máquina de destino. & Obligatorio. 36 car. \\
\texttt{description} & STRING & Justificación técnica del vínculo. & Obligatorio. 10 a 1000 car. \\ \bottomrule
\end{tabular}
\end{table}

\textbf{Ejemplo de Relación:}
\begin{lstlisting}[language=json]
{
  "id": "111e8400-e29b-41d4-a716-446655440111",
  "source": "LeadSynth",
  "destination": "MainMixer",
  "description": "Routes audio from lead to master mixer."
}
\end{lstlisting}

\section{Reglas de Validación}

Para garantizar la integridad del flujo entre máquinas, se aplican las siguientes restricciones:

\subsection{Identidad y Formato}
\begin{itemize}
    \item \textbf{Unicidad de IDs:} Cada relación en el array \texttt{relations} debe tener un \texttt{id} único.
    \item \textbf{Formato UUID:} El campo \texttt{id} debe cumplir con el estándar de 36 caracteres.
\end{itemize}

\subsection{Restricciones de Texto}
\begin{itemize}
    \item \textbf{Descripción Global:} Máximo 1000 caracteres.
    \item \textbf{Descripción de Relación:} Entre 10 y 1000 caracteres.
\end{itemize}

\subsection{Interfaces Expuestas y Existencia}
\begin{itemize}
    \item \textbf{Existencia:} Los IDs especificados en \texttt{source} y \texttt{destination} deben corresponder a máquinas CIM existentes en el proyecto.
    \item \textbf{Sin Auto-Relación:} Una máquina no puede conectarse consigo misma. Los campos \texttt{source} y \texttt{destination} deben ser diferentes.
    \item \textbf{Configuración de Origen (\texttt{source}):} La máquina referenciada como \textbf{origen} DEBE contener al menos un elemento interno configurado con \texttt{hasExternalOutput: true}.
    \item \textbf{Configuración de Destino (\texttt{destination}):} La máquina referenciada como \textbf{destino} DEBE contener al menos un elemento interno configurado con \texttt{hasExternalInput: true}.
\end{itemize}

En resumen, las relaciones CIM solo pueden existir si las interfaces individuales de las máquinas están habilitadas para la comunicación externa y si ambas entidades son máquinas distintas y existentes.

\section{Guía de Uso}

Las relaciones CIM expresan dependencias conceptuales (ej.: \textit{La máquina de síntesis de leads alimenta a la máquina de mezcla principal}). No entran en detalles técnicos de canales, puertos específicos o parámetros muy concretos, ya que esa granularidad se define exclusivamente en el nivel \textbf{PIM}.

\begin{tcolorbox}[colback=blue!5!white,colframe=blue!75!black,title=CONSEJO]
Antes de definir una relación, asegúrese de que la máquina origen expone salidas y la destino acepta entradas en sus respectivos modelos CIM. De lo contrario, el validador bloqueará la creación del flujo.
\end{tcolorbox}
